In der vorangegangenen Aufgabe wurde die ganze Funktion $f(z)=(e^z-1)/z$
in eine Potenzreihe entwickelt.
\begin{teilaufgaben}
\item Bestimmen Sie die Nullstellen der Funktion $f(z)$.
\item Die reziproke Funktion hat die Potenzreihenentwicklung
\begin{equation}
g(z)
=
\frac{1}{f(z)}
=
\frac{z}{e^z-1}
=
\sum_{n=0}^\infty \frac{B_n}{n!}z^k.
\label{buch:402:greihe}
\end{equation}
Die Koeffizienten $B_n$ heissen die Bernoulli-Zahlen.
Bestimmen Sie die ersten Bernoulli-Zahlen bis $B_4$, indem 
Sie den Ansatz~\eqref{buch:402:greihe}
in die Gleichung
\[
z = (e^z-1) g(z)
\]
einsetzen und die Koeffizienten durch Koeffizientenvergleich
bestimmen.
\item
Bestimmen Sie den Konvergenzradius der Reihe.
\end{teilaufgaben}

\begin{loesung}
\begin{teilaufgaben}
\item
Die Nullstellen von $f(z)$ sind die Nullstellen von $e^z-1$ oder die
Lösungen der Gleichung
\begin{equation}
e^a(\cos b + i \sin b)
=
1.
\label{buch:402:gleichung}
\end{equation}
Daraus liest man ab, dass $\sin b=0$ sein muss,
d.~h. $b\in \pi\mathbb{Z}$.
Für die ungeraden Vielfachen von $\pi$ führen zu einem negativen $\cos b$,
es bleiben daher nur die geraden Vielfachen von $\pi$.
Um die Gleichung~\eqref{buch:402:gleichung} zu erfüllen muss jetzt nur noch
$e^a=1$ sein, d.~h.~$a=0$.
Die Menge der Nullstellen ist daher $2\pi i\mathbb{Z}\setminus\{0\}$.
\item
Wir schreiben Terme bis vierter Ordnung der beiden Reihen und
multiplizieren aus
\begin{align*}
z
&=
\biggl(
\frac{z}{2} + \frac{z^2}{3!} + \frac{z^3}{4!} + \frac{z^4}{5!} + \ldots
\biggr)
\biggl(
B_0 + B_1z + \frac{B_2}{2!}z^2 + \frac{B_3}{3!}z^3 + \frac{B_4}{4!}z^4 + \ldots
\biggr)
\\
&=
\textstyle
B_0z
+
\bigl(\frac12 B_0 + B_1\bigr) z^2
+
\bigl(\frac16 B_0 + \frac12 B_1 + \frac12 B_2\bigr)z^2
+
\bigl(\frac1{24} B_0 + \frac16 B_1 + \frac14 B_2 + \frac16 B_3\bigr)z^3
\\
&\qquad
\textstyle
+
\bigl(\frac1{120} B_0 + \frac1{24} B_1 + \frac1{12} B_2 + \frac1{12} B_3 + \frac1{24} B_4)z^4
+
\ldots
\end{align*}
Daraus kann man die Gleichungen
\[
\renewcommand{\arraystretch}{1.4}
\renewcommand{\arraycolsep}{2pt}
\begin{array}{rcrcrcrcrcl}
           B_0 & &               & &               & &               & &               &=& 1\phantom{.}\\
\frac1{2}  B_0 &+&           B_1 & &               & &               & &               &=& 0\phantom{.}\\
\frac1{6}  B_0 &+& \frac1{2} B_1 &+& \frac1{2} B_2 & &               & &               &=& 0\phantom{.}\\
\frac1{24} B_0 &+& \frac1{6} B_1 &+& \frac1{4} B_2 &+& \frac1{6} B_3 & &               &=& 0\phantom{.}\\
\frac1{120}B_0 &+& \frac1{24}B_1 &+& \frac1{12}B_2 &+& \frac1{12}B_3 &+& \frac1{24}B_4 &=& 0.
\end{array}
\]
Da das Gleichungssystem Dreiecksform hat, kann man die
Koeffizienten sofort rekursiv auflösen und findet
\[
\renewcommand{\arraystretch}{1.4}
\renewcommand{\arraycolsep}{2pt}
\begin{array}{rclcr}
B_0 &=& &=&\phantom{-}1
\\
B_1 &=& &=&-\frac12
\\
B_2 &=& -\frac{1}{3} B_0 - B_1 = -\frac13 + \frac12 &=& \phantom{-}\frac{1}{6}
\\
B_3 &=& -\frac1{4}   B_0 - B_1 - \frac{3}{2} B_2
     = -\frac14 + \frac12 - \frac14
     &=& \phantom{-}0
\\
B_4 &=& -\frac1{5}   B_0 - B_1 - 2           B_2 - 2 B_3
     = -\frac1{5} +\frac12 -\frac13
     &=& -\frac1{30}\rlap{.}
\end{array}
\]
Damit sind die ersten Bernoulli-Zahlen bis $B_4$ bestimmt, wie in der
Aufgabe verlangt.

Man kann mit der gleichen Methode auch mehr erreichen, als in der
Aufgabe verlangt war.
Man kann nämlich allgemein
\begin{align*}
z
&=
\biggl(
\sum_{k=0}^\infty \frac{z^k}{(k+1)!}
\biggr)
\biggl(
\sum_{n=0}^\infty \frac{B_n}{n!}z^n
\biggr)
=
\sum_{l=0}^\infty
\biggl(
\sum_{n=0}^l
\frac{B_n}{n!}
\cdot
\frac{1}{(l-n+1)!}
\biggr)
z^l
\end{align*}
schreiben.
Für $l>1$ kann man daraus die Rekursionsformel
\[
\frac{B_l}{l!}
+
\sum_{n=0}^{l-1}
\frac{B_n}{n!}
\cdot
\frac{1}{(l-n+1)!}
=
0
\qquad\Rightarrow\qquad
\frac{B_l}{l!}
=
-
\sum_{n=0}^{l-1}
\frac{B_n}{n!}
\cdot
\frac{1}{(l-n+1)!}
\]
ableiten.
Tatsächlich gibt die Formel für $l=4$
\begin{align*}
0
&=
\frac{B_4}{4!}
+
\frac{B_3}{3!}\frac{1}{2!}
+
\frac{B_2}{2!}\frac{1}{3!}
+
\frac{B_1}{1!}\frac{1}{4!}
+
\frac{B_0}{0!}\frac{1}{5!}
\\
&=
\frac{1}{24}B_4
+
\frac{1}{12}B_3
+
\frac{1}{12}B_2
+
\frac{1}{24}B_1
+
\frac{1}{120}B_0,
\end{align*}
was mit der oben gefundenen Formel übereinstimmt.
Multipliziert man mit $l!$ folgt
\begin{align*}
B_l
&=
-
\sum_{n=0}^{l-1}
\frac{l!}{n!\,(l-n+1)!} B_n
=
-
\sum_{n=0}^{l-1}
\frac{1}{l-n+1}
\binom{l}{n}
B_n.
\end{align*}
\item
Die Funktion $g(z)$ hat Pole bei den Nullstellen von $f(z)$.
Die Pole sind $2\pi ik$, $k\in \mathbb{Z}\setminus\{0\}$, daher ist
der nächste Pol zum Nullpunkt der Punkt $z=\pm 2\pi i$ und daher 
ist der Konvergenzradius $=2\pi$.
\qedhere
\end{teilaufgaben}
\end{loesung}

